\documentclass[10pt]{article}

\usepackage[margin=0.75in]{geometry}
\usepackage{amsmath,amsthm,amssymb,bm}
\usepackage{xcolor}
\usepackage{cancel}
\usepackage{graphicx}
\usepackage{changepage}
\usepackage{circuitikz}
\usepackage{pgfplots}
\usepackage{physics}
\usepackage{hyperref}
\usepackage{siunitx}
\usepackage{minted}
\usepackage{fontspec}
\usepackage{relsize}
\usepackage{subfig}
\usepackage{mhchem}
\usepackage{todonotes}
\usepackage{biblatex}
\usepackage{multicol, multirow, booktabs}
\usepackage[breakable]{tcolorbox}
\usepackage[inline]{enumitem}
\patchcmd{\thebibliography}{\section*{\refname}}{}{}{}
\addbibresource{sources.bib}

\theoremstyle{definition}
\newtheorem{problem}{Problem}
\newtheorem{soln}{Solution}

\pgfplotsset{compat=newest}
\usetikzlibrary{lindenmayersystems}
\usetikzlibrary{arrows}
\usetikzlibrary{calc}
\usetikzlibrary{positioning, fit}
\usetikzlibrary{3d, perspective}

\definecolor{incolor}{HTML}{303F9F}
\definecolor{outcolor}{HTML}{D84315}
\definecolor{cellborder}{HTML}{CFCFCF}
\definecolor{cellbackground}{HTML}{F7F7F7}
\newcommand{\ui}{\hat{i}}
\newcommand{\uj}{\hat{j}}
\newcommand{\uk}{\hat{k}}
\newcommand{\ux}{\hat{x}}
\newcommand{\uy}{\hat{y}}
\newcommand{\uz}{\hat{z}}
\newcommand{\uv}[1]{\hat{\bm{{#1}}}}
\newcommand{\pr}[1]{{#1^\prime}}
\newcommand{\justif}[2]{&{#1}&\text{#2}}
\newcommand{\dul}[1]{\underline{\underline{#1}}}
\newcommand{\pdiff}[2][]{{\frac{\partial^{#1}}{\partial {{#2}^{#1}}}}}
% \newcommand{\dif}[2][]{{\frac{d^{#1}}{d{{#2}^{#1}}}}}
\pgfdeclarelayer{background}  
\pgfsetlayers{background,main}
\AtBeginDocument{\RenewCommandCopy\qty\SI}

\makeatletter
\newcommand{\boxspacing}{\kern\kvtcb@left@rule\kern\kvtcb@boxsep}
\makeatother
\newcommand{\prompt}[4]{
    \ttfamily\llap{{\color{#2}[#3]:\hspace{3pt}#4}}\vspace{-\baselineskip}
}

\newcommand{\thevenin}[2]{
  \begin{center}
    \begin{circuitikz} \draw
      (0,0) -- (2,0) to[battery1, l_=$V_{Th}\eq#1$] (2,2) 
      to[resistor, l_=$R_{Th}\eq#2$] (0,2)
      ;
      \draw [o-] (-.07,2.079);
      \draw [o-] (-.07,0.079);
    \end{circuitikz}
  \end{center}
}

\newcommand{\norton}[2]{
  \begin{center}
    \begin{circuitikz} \draw
      (0,0) -- (3,0) to[american current source, l_=$I_{N}\eq#1$] (3,2) -- (0,2) (2,0)
      to[resistor, l=$R_{N}\eq#2$] (2,2)
      ;
      \draw [o-] (-.07,2.079);
      \draw [o-] (-.07,0.079);
    \end{circuitikz}
  \end{center}
}

\newcommand{\highlight}[1]{\colorbox{yellow}{#1}}

\newcommand{\ti}[1]{\widetilde{#1}}

\newfontface{\Kaufmann}{Kaufmann}
\DeclareTextFontCommand{\kf}{\Kaufmann}
\newcommand{\scriptr}{\fontsize{12pt}{12pt}\kf{r}}

\newfontface{\KaufmannB}{Kaufmann Bold}
\DeclareTextFontCommand{\kfb}{\KaufmannB}
\newcommand{\bscriptr}{\fontsize{12pt}{12pt}\kfb{r}}

\newcommand{\bv}[1]{\vec{#1}}

\title{Physics 4605H: Assignment VII}
\author{Jeremy Favro (0805980) \\ Trent University, Peterborough, ON, Canada}
\date{\today}

\begin{document}
\maketitle

% PROBLEM 1
\begin{problem}
Construct the two-qubit $\text{CNOT}_{21}$ gate.
The subscript $21$ here means that qubit $2$ is the control and qubit $1$ is the target, where
qubit 1 is defined to be the qubit which appears first in the notation $\ket{00}$ and qubit 2 is the
one which appears second.
Explain what you are doing, and show the gate both as a matrix and in Dirac/operator
notation.
\end{problem}
\begin{soln}
  The way I was taught to do this in 2605 (and one of the ways we did it in class) with Dirac notation was to write my desired
  outcomes per qubit (superscript $(x)$ denotes things in the qubit $x$ Hilbert space):
  \begin{enumerate}[label=(\roman*)]
    \item If qubit 2 is zero, do nothing to the state of qubit 1 (apply the identity operator to it),
          $$\mathbb{I}^{(1)}\otimes \ket{0}\bra{0}^{(2)}$$
    \item If qubit 2 is 1, flip qubit 1's state (apply $\sigma_x$ to it),
          $$\hat{\sigma}_x^{(1)}\otimes \ket{1}\bra{1}^{(2)}.$$
  \end{enumerate}
  Then sum these to create a two-condition operator. Note that order matters and qubit operators should be tensored
  in ascending order of qubit index.
  In matrix notation,
  $$\mathbb{I}^{(1)}\otimes \ket{0}\bra{0}^{(2)}=\begin{pmatrix}
      1 & 0 \\
      0 & 1
    \end{pmatrix}\otimes
    \begin{pmatrix}
      1 & 0 \\
      0 & 0
    \end{pmatrix}
    =\begin{pmatrix}
      1 & 0 & 0 & 0 \\
      0 & 0 & 0 & 0 \\
      0 & 0 & 1 & 0 \\
      0 & 0 & 0 & 0
    \end{pmatrix}
  $$
  and
  $$\hat{\sigma}_x^{(1)}\otimes \ket{1}\bra{1}^{(2)}=\begin{pmatrix}
      0 & 1 \\
      1 & 0
    \end{pmatrix}\otimes \begin{pmatrix}
      0 & 0 \\
      0 & 1
    \end{pmatrix}
    =\begin{pmatrix}
      0 & 0 & 0 & 0 \\
      0 & 0 & 0 & 1 \\
      0 & 0 & 0 & 0 \\
      0 & 1 & 0 & 0
    \end{pmatrix}
  $$
  and their sum is
  $$\begin{pmatrix}
      1 & 0 & 0 & 0 \\
      0 & 0 & 0 & 0 \\
      0 & 0 & 1 & 0 \\
      0 & 0 & 0 & 0
    \end{pmatrix}+\begin{pmatrix}
      0 & 0 & 0 & 0 \\
      0 & 0 & 0 & 1 \\
      0 & 0 & 0 & 0 \\
      0 & 1 & 0 & 0
    \end{pmatrix}
    =\begin{pmatrix}
      1 & 0 & 0 & 0 \\
      0 & 0 & 0 & 1 \\
      0 & 0 & 1 & 0 \\
      0 & 1 & 0 & 0
    \end{pmatrix}.
  $$
  This is, to be explicit,
  $$\mathbb{I}^{(1)}\otimes \ket{0}\bra{0}^{(2)} +\hat{\sigma}_x^{(1)}\otimes \ket{1}\bra{1}^{(2)}$$
  in Dirac notation. After a discussion in office hours I will also show how we can get this from a truth table.
  Recall that we can write a matrix for an operator $\hat{A}$ acting on a two-qubit system as
  $$\begin{pmatrix}
      \bra{00}\hat{A}\ket{00} & \bra{00}\hat{A}\ket{01} & \bra{00}\hat{A}\ket{10} & \bra{00}\hat{A}\ket{11} \\
      \bra{01}\hat{A}\ket{00} & \dots                   &                         &                         \\
      \bra{10}\hat{A}\ket{00} & \dots                   &                         &                         \\
      \bra{11}\hat{A}\ket{00} & \dots                   &                         &                         \\
    \end{pmatrix}
    .$$
  We also know that the operator that we want has the following truth table
  $$
    \begin{array}{|c|c|}
      \ket{\psi} & \hat{A}\ket{\psi} \\
      \hline
      \ket{00}   & \ket{00}          \\
      \ket{01}   & \ket{11}          \\
      \ket{10}   & \ket{10}          \\
      \ket{11}   & \ket{01}
    \end{array}
  $$
  and hence we can evaluate the entries of the matrix we defined. Note that we can speed the process up a bit by representing the matrix like
  $$\begin{pmatrix}
      \bra{00} & \bra{11} & \bra{10} & \bra{01} \\
      \bra{01} &          &          &          \\
      \bra{10} &          &          &          \\
      \bra{11} &          &          &          \\
    \end{pmatrix}
  $$
  and noting that when we ``operate'' on this with $\hat{A}$ as defined by the truth table we obtain the following using the representation
  $$
    \begin{pmatrix}
      \bra{00} & \bra{00} & \bra{11} & \bra{10} \\
      \bra{01} &          &          &          \\
      \bra{10} &          &          &          \\
      \bra{11} &          &          &          \\
    \end{pmatrix}
  $$
  and from here it is easy to see that when we compute the inner products we obtain the same matrix as with the other method,
  $$
    \begin{pmatrix}
      1 & 0 & 0 & 0 \\
      0 & 0 & 0 & 1 \\
      0 & 0 & 1 & 0 \\
      0 & 1 & 0 & 0
    \end{pmatrix}
  $$
\end{soln}

% PROBLEM 2
\begin{problem}
We discussed in class how to implement a $\text{CNOT}_{12}$ gate using J coupling in NMR. We
concluded that
$$\hat{R}_z^{(1)}(\pi/2)\hat{R}_y^{(2)}(-\pi/2)\hat{R}_z^{(2)}(-\pi/2)\hat{U}_J^{(2)}(1/2J_{12})\hat{R}_y^{(2)}(\pi/2)=e^{i\pi/4}\hat{U}_{CNOT}$$
where
$$\hat{U}_J^{(2)}(1/2J_{12})\equiv \ket{0}\bra{0}^{(1)}e^{-i(\pi/4)\hat{\sigma}_z^{(2)}}+\ket{1}\bra{1}^{(1)}e^{i(\pi/4)\hat{\sigma}_z^{(2)}}$$
The goal of this problem is to demonstrate that this is true by writing out each of the rotation
operators in terms of $\hat{\sigma}_i$ operators and multiplying the terms. Follow these steps:
\begin{enumerate}[label=(\alph*)]
  \item Write out each of the rotation operators.
        For example, $\hat{R}_y^{(2)}(\pi/2)=\cos(\pi/4)\mathbb{I}^{(2)}-i\sin(\pi/4)\hat{\sigma}^{(2)}_y=(\mathbb{I}^{(2)}-i\hat{\sigma}_y^{(2)})/\sqrt{2}$.
  \item Calculate $\hat{U}_J^{(2)}(1/2J_{12})\hat{R}_y^{(2)}(\pi/2)$. Simplify such that each term contains no more than
        one $\hat{\sigma}_i$ operator. Keep in mind that $\hat{\sigma}_i\hat{\sigma}_j = \delta_{ij} + i\hat{\sigma}_k \epsilon _{ijk}$.
  \item Calculate $\hat{R}_z^{(2)}(-\pi/2)\hat{U}_J^{(2)}(1/2J_{12})\hat{R}_y^{(2)}(\pi/2)$. Simplify again.
  \item Calculate $\hat{R}_y^{(2)}(-\pi/2)\hat{R}_z^{(2)}(-\pi/2)\hat{U}_J^{(2)}(1/2J_{12})\hat{R}_y^{(2)}(\pi/2)$. Simplify again.
  \item Calculate $\hat{R}_z^{(1)}(\pi/2)\hat{R}_y^{(2)}(-\pi/2)\hat{R}_z^{(2)}(-\pi/2)\hat{U}_J^{(2)}(1/2J_{12})\hat{R}_y^{(2)}(\pi/2)$. Simplify and
        demonstrate that your result equals $e^{i\pi/4}\hat{U}_{CNOT}$.
\end{enumerate}
\end{problem}
\begin{soln}
  \begin{enumerate}[label=(\alph*)]
    \item Recall that we defined the rotation operator as
          $$\hat{R}_{\uv{j}}=\exp\left(-i\frac{\theta}{2}\uv{j}\cdot\vec{\hat{\sigma}}\right)=\cos\left(\frac{\theta}{2}\right)\hat{I}-i\sin\left(\frac{\theta}{2}\right)(\uv{j}\cdot \vec{\hat{\sigma}})$$
          where $\uv{j}$ is the unit vector for the axis we want to rotate about. The operators we want are then:
          \begin{enumerate}[label=(\roman*)]
            \item $\hat{R}_z^{(1)}(\pi/2)=\exp\left(-i\frac{\pi}{4}\hat{\sigma}_z^{(1)}\right)
                    =\cos\left(\frac{\pi}{4}\right)\hat{I}^{(1)}-i\sin\left(\frac{\pi}{4}\right)\hat{\sigma}_z^{(1)}
                    =(\hat{I}^{(1)}-i\hat{\sigma}_z^{(1)})/\sqrt{2}$
            \item $\hat{R}_y^{(2)}(-\pi/2)=\exp\left(i\frac{\pi}{4}\hat{\sigma}_y^{(2)}\right)
                    =\cos\left(\frac{\pi}{4}\right)\hat{I}^{(2)}+i\sin\left(\frac{\pi}{4}\right)\hat{\sigma}_y^{(2)}
                    =(\hat{I}^{(2)}+i\hat{\sigma}_y^{(2)})/\sqrt{2}
                  $
            \item $\hat{R}_z^{(2)}(-\pi/2)=\exp\left(i\frac{\pi}{4}\hat{\sigma}_z^{(2)}\right)
                    =\cos\left(\frac{\pi}{4}\right)\hat{I}^{(2)}+i\sin\left(\frac{\pi}{4}\right)\hat{\sigma}_z^{(2)}
                    =(\hat{I}^{(2)}+i\hat{\sigma}_z^{(2)})/\sqrt{2}
                  $
            \item $
                    \hat{R}_y^{(2)}(\pi/2)
                    =\exp\left(-i\frac{\pi}{4}\hat{\sigma}_y^{(2)}\right)
                    =\cos\left(\frac{\pi}{4}\right)\hat{I}^{(2)}-i\sin\left(\frac{\pi}{4}\right)\hat{\sigma}_y^{(2)}
                    =(\hat{I}^{(2)}-i\hat{\sigma}_y^{(2)})/\sqrt{2}
                  $
          \end{enumerate}
    \item I'm going to do this in two parts per step, one for the $\ket{0}\bra{0}$ terms and another for the $\ket{1}\bra{1}$ terms, then I'll recombine them at the end:
          \begin{enumerate}[label=(\arabic*)]
            \item Note that
                  $$\ket{0}\bra{0}^{(1)}e^{-i(\pi/4)\hat{\sigma}_z^{(2)}}=\ket{0}\bra{0}^{(1)}e^{-i(\pi/4)\hat{\sigma}_z^{(2)}}=\ket{0}\bra{0}^{(1)}\otimes(\hat{I}^{(2)}-i\hat{\sigma}_z^{(2)})/\sqrt{2}.$$
                  We can pull it apart this way by rewriting the exponential using Euler's formula
                  and writing the sine and cosine as their respective power series which we can then simplify as the cosine series is all even powers so the $\sigma_z$s becomes the identity
                  and the sine is all odd so we pick up a factor of $\sigma_z$. This works for the other Pauli matrices as well as they all have the square-identity property. So:
                  \begin{align*}
                    \ket{0}\bra{0}^{(1)}\otimes(\hat{I}^{(2)}-i\hat{\sigma}_z^{(2)})/\sqrt{2}\hat{R}_y^{(2)}(\pi/2) & =\ket{0}\bra{0}^{(1)}\otimes(\hat{I}^{(2)}-i\hat{\sigma}_z^{(2)})/\sqrt{2}(\hat{I}^{(2)}-i\hat{\sigma}_y^{(2)})/\sqrt{2}                                                                \\
                                                                                                                    & =\ket{0}\bra{0}^{(1)}\otimes(\hat{I}^{(2)}-i\hat{\sigma}_z^{(2)})(\hat{I}^{(2)}-i\hat{\sigma}_y^{(2)})/2                                                                                \\
                                                                                                                    & =\frac{1}{2}\ket{0}\bra{0}^{(1)}\otimes((\hat{I}^{(2)}\hat{I}^{(2)}-i\hat{I}^{(2)}\hat{\sigma}_y^{(2)})-i(\hat{\sigma}_z^{(2)}\hat{I}^{(2)}-i\hat{\sigma}_z^{(2)}\hat{\sigma}_y^{(2)})) \\
                                                                                                                    & =\frac{1}{2}\ket{0}\bra{0}^{(1)}\otimes(\hat{I}^{(2)}-i\hat{\sigma}_y^{(2)}-i\hat{\sigma}_z^{(2)}-\hat{\sigma}_z^{(2)}\hat{\sigma}_y^{(2)})                                             \\
                  \end{align*}
                  and using the Levi-Cevita notation
                  $$\hat{\sigma}_z^{(2)}\hat{\sigma}_y^{(2)}=-i\hat{\sigma}_x^{(2)}$$
                  so
                  \begin{align*}
                    \ket{0}\bra{0}^{(1)}\otimes(\hat{I}^{(2)}-i\hat{\sigma}_z^{(2)})/\sqrt{2}\hat{R}_y^{(2)}(\pi/2) & =\frac{1}{2}\ket{0}\bra{0}^{(1)}\otimes(\hat{I}^{(2)}-i\hat{\sigma}_y^{(2)}-i\hat{\sigma}_z^{(2)}+i\hat{\sigma}_x^{(2)}) \\
                  \end{align*}
            \item \begin{align*}
                    \ket{1}\bra{1}^{(1)}e^{i(\pi/4)\hat{\sigma}_z^{(2)}}\hat{R}_y^{(2)}(\pi/2) & =\frac{1}{2}\ket{1}\bra{1}^{(1)}\otimes(\hat{I}^{(2)}+i\hat{\sigma}_z^{(2)})(\hat{I}^{(2)}-i\hat{\sigma}_y^{(2)})                                                                  \\
                                                                                               & =\frac{1}{2}\ket{1}\bra{1}^{(1)}\otimes(\hat{I}^{(2)}+i\hat{\sigma}_z^{(2)})(\hat{I}^{(2)}-i\hat{\sigma}_y^{(2)})                                                                  \\
                                                                                               & =\frac{1}{2}\ket{1}\bra{1}^{(1)}\otimes(\hat{I}^{(2)}(\hat{I}^{(2)}-i\hat{\sigma}_y^{(2)})+i\hat{\sigma}_z^{(2)}(\hat{I}^{(2)}-i\hat{\sigma}_y^{(2)}))                             \\
                                                                                               & =\frac{1}{2}\ket{1}\bra{1}^{(1)}\otimes(\hat{I}^{(2)}\hat{I}^{(2)}-i\hat{I}^{(2)}\hat{\sigma}_y^{(2)}+i\hat{\sigma}_z^{(2)}\hat{I}^{(2)}+\hat{\sigma}_z^{(2)}\hat{\sigma}_y^{(2)}) \\
                                                                                               & =\frac{1}{2}\ket{1}\bra{1}^{(1)}\otimes(\hat{I}^{(2)}-i\hat{\sigma}_y^{(2)}+i\hat{\sigma}_z^{(2)}-i\hat{\sigma}_x^{(2)})                                                           \\
                  \end{align*}
          \end{enumerate}
    \item ~\begin{enumerate}[label=(\arabic*)]
            \item \begin{align*}
                    \hat{R}_z^{(2)}(-\pi/2)\hat{U}_J^{(2)}(1/2J_{12})\hat{R}_y^{(2)}(\pi/2) & =_{(0)} (\hat{I}^{(2)}+i\hat{\sigma}_z^{(2)})\frac{1}{2\sqrt{2}}\ket{0}\bra{0}^{(1)}\otimes(\hat{I}^{(2)}-i\hat{\sigma}_y^{(2)}-i\hat{\sigma}_z^{(2)}+i\hat{\sigma}_x^{(2)}) \\
                                                                                            & = \frac{1}{2\sqrt{2}}(\hat{I}^{(2)}+i\hat{\sigma}_z^{(2)})\ket{0}\bra{0}^{(1)}\otimes(\hat{I}^{(2)}-i\hat{\sigma}_y^{(2)}-i\hat{\sigma}_z^{(2)}+i\hat{\sigma}_x^{(2)})       \\
                  \end{align*}
            \item \begin{align*}
                    \hat{R}_z^{(2)}(-\pi/2)\hat{U}_J^{(2)}(1/2J_{12})\hat{R}_y^{(2)}(\pi/2) & =_{(1)}\frac{1}{2\sqrt{2}}(\hat{I}^{(2)}+i\hat{\sigma}_z^{(2)})\ket{1}\bra{1}^{(1)}\otimes(\hat{I}^{(2)}-i\hat{\sigma}_y^{(2)}+i\hat{\sigma}_z^{(2)}-i\hat{\sigma}_x^{(2)}) \\
                  \end{align*}
          \end{enumerate}
    \item  ~\begin{enumerate}[label=(\arabic*)]
            \item \begin{align*}
                    \hat{R}_y^{(2)}(-\pi/2)\hat{R}_z^{(2)}(-\pi/2)\hat{U}_J^{(2)}(1/2J_{12})\hat{R}_y^{(2)}(\pi/2) & = \frac{1}{4}(\hat{I}^{(2)}+i\hat{\sigma}_y^{(2)})(\hat{I}^{(2)}+i\hat{\sigma}_z^{(2)})\ket{0}\bra{0}^{(1)}\otimes(\hat{I}^{(2)}-i\hat{\sigma}_y^{(2)}-i\hat{\sigma}_z^{(2)}+i\hat{\sigma}_x^{(2)}) \\
                                                                                                                   & = \frac{1}{4}(\hat{I}^{(2)}\hat{I}^{(2)}+i\hat{I}^{(2)}\hat{\sigma}_z^{(2)}+i\hat{\sigma}_y^{(2)}\hat{I}^{(2)}-\hat{\sigma}_y^{(2)}\hat{\sigma}_z^{(2)})                                            \\
                                                                                                                   & \quad\ket{0}\bra{0}^{(1)}\otimes(\hat{I}^{(2)}-i\hat{\sigma}_y^{(2)}-i\hat{\sigma}_z^{(2)}+i\hat{\sigma}_x^{(2)})                                                                                   \\
                                                                                                                   & = \frac{1}{4}(\hat{I}^{(2)}+i\hat{\sigma}_z^{(2)}+i\hat{\sigma}_y^{(2)}-\hat{\sigma}_x^{(2)})                                                                                                       \\
                                                                                                                   & \quad\ket{0}\bra{0}^{(1)}\otimes(\hat{I}^{(2)}-i\hat{\sigma}_y^{(2)}-i\hat{\sigma}_z^{(2)}+i\hat{\sigma}_x^{(2)})                                                                                   \\
                  \end{align*}
            \item \begin{align*}
                    \hat{R}_y^{(2)}(-\pi/2)\hat{R}_z^{(2)}(-\pi/2)\hat{U}_J^{(2)}(1/2J_{12})\hat{R}_y^{(2)}(\pi/2) & =_{(1)}\frac{1}{4}(\hat{I}^{(2)}+i\hat{\sigma}_y^{(2)})(\hat{I}^{(2)}+i\hat{\sigma}_z^{(2)})\ket{1}\bra{1}^{(1)}\otimes(\hat{I}^{(2)}-i\hat{\sigma}_y^{(2)}+i\hat{\sigma}_z^{(2)}-i\hat{\sigma}_x^{(2)}) \\
                                                                                                                   & =_{(1)}\frac{1}{4}(\hat{I}^{(2)}+i\hat{\sigma}_z^{(2)}+i\hat{\sigma}_y^{(2)}-\hat{\sigma}_x^{(2)})                                                                                                       \\
                                                                                                                   & \quad\ket{1}\bra{1}^{(1)}\otimes(\hat{I}^{(2)}-i\hat{\sigma}_y^{(2)}+i\hat{\sigma}_z^{(2)}-i\hat{\sigma}_x^{(2)})                                                                                        \\
                  \end{align*}
          \end{enumerate}
          \newpage
    \item ~\begin{enumerate}[label=(\arabic*)]
            \item With $\hat{R}_z^{(1)}(\pi/2)\hat{R}_y^{(2)}(-\pi/2)\hat{R}_z^{(2)}(-\pi/2)\hat{U}_J^{(2)}(1/2J_{12})\hat{R}_y^{(2)}(\pi/2)=\hat{A}$,

                  \begin{align*}
                    \hat{A} & = \frac{1}{4\sqrt{2}}(\hat{I}^{(1)}-i\hat{\sigma}_z^{(1)})(\hat{I}^{(2)}+i\hat{\sigma}_z^{(2)}+i\hat{\sigma}_y^{(2)}-\hat{\sigma}_x^{(2)}) \\
                            & \quad\ket{0}\bra{0}^{(1)}\otimes(\hat{I}^{(2)}-i\hat{\sigma}_y^{(2)}-i\hat{\sigma}_z^{(2)}+i\hat{\sigma}_x^{(2)})                          \\
                            & = \frac{1}{4\sqrt{2}}
                    (
                    \hat{I}^{(1)}\hat{I}^{(2)}+i\hat{I}^{(1)}\hat{\sigma}_z^{(2)}+i\hat{I}^{(1)}\hat{\sigma}_y^{(2)}-\hat{I}^{(1)}\hat{\sigma}_x^{(2)}
                    +\hat{\sigma}_z^{(1)}\hat{I}^{(2)}+\hat{\sigma}_z^{(1)}\hat{\sigma}_z^{(2)}+\hat{\sigma}_z^{(1)}\hat{\sigma}_y^{(2)}+i\hat{\sigma}_z^{(1)}\hat{\sigma}_x^{(2)}
                    )
                    \\
                            & \quad\ket{0}\bra{0}^{(1)}\otimes(\hat{I}^{(2)}-i\hat{\sigma}_y^{(2)}-i\hat{\sigma}_z^{(2)}+i\hat{\sigma}_x^{(2)})                          \\
                  \end{align*}
            \item \begin{align*}
                    \hat{A} &
                    =_{(1)}\frac{1}{4\sqrt{2}}
                    (
                    \hat{I}^{(1)}\hat{I}^{(2)}+i\hat{I}^{(1)}\hat{\sigma}_z^{(2)}+i\hat{I}^{(1)}\hat{\sigma}_y^{(2)}-\hat{I}^{(1)}\hat{\sigma}_x^{(2)}
                    +\hat{\sigma}_z^{(1)}\hat{I}^{(2)}+\hat{\sigma}_z^{(1)}\hat{\sigma}_z^{(2)}+\hat{\sigma}_z^{(1)}\hat{\sigma}_y^{(2)}+i\hat{\sigma}_z^{(1)}\hat{\sigma}_x^{(2)}
                    )
                    \\
                            & \quad\ket{1}\bra{1}^{(1)}\otimes(\hat{I}^{(2)}-i\hat{\sigma}_y^{(2)}+i\hat{\sigma}_z^{(2)}-i\hat{\sigma}_x^{(2)}) \\
                  \end{align*}
          \end{enumerate}
          so recombining the two terms the whole operator is
          \begin{align*}
             & \frac{1}{4\sqrt{2}}
            (
            \hat{I}^{(1)}\hat{I}^{(2)}+i\hat{I}^{(1)}\hat{\sigma}_z^{(2)}+i\hat{I}^{(1)}\hat{\sigma}_y^{(2)}-\hat{I}^{(1)}\hat{\sigma}_x^{(2)}
            +\hat{\sigma}_z^{(1)}\hat{I}^{(2)}+\hat{\sigma}_z^{(1)}\hat{\sigma}_z^{(2)}+\hat{\sigma}_z^{(1)}\hat{\sigma}_y^{(2)}+i\hat{\sigma}_z^{(1)}\hat{\sigma}_x^{(2)}
            )                                                                                                                                                                                                                                          \\
             & \left[\ket{0}\bra{0}^{(1)}\otimes(\hat{I}^{(2)}-i\hat{\sigma}_y^{(2)}-i\hat{\sigma}_z^{(2)}+i\hat{\sigma}_x^{(2)})+ \ket{1}\bra{1}^{(1)}\otimes(\hat{I}^{(2)}-i\hat{\sigma}_y^{(2)}+i\hat{\sigma}_z^{(2)}-i\hat{\sigma}_x^{(2)})\right]
          \end{align*}
          This is evidently not the $\text{CNOT}_{12}$ gate unless I'm missing something. I can't tell where I went wrong.


          % \item It's easiest (for me) to do this with the exponential forms we calculated previously. Note that when we are applying an operator like
          %       $$\hat{R}_y^{(2)}(\pi/2)
          %         =\exp\left(-i\frac{\pi}{4}\hat{\sigma}_y^{(2)}\right)$$
          %       to a two qubit system, what we are really doing is applying
          %       $$\hat{R}_y^{(2)}(\pi/2)
          %         =\exp\left(-i\frac{\pi}{4}\hat{I}^{(1)}\otimes\hat{\sigma}_y^{(2)}\right)$$
          %       i.e. we do nothing to the first qubit, and rotate the second. Also recall that when we original found $\hat{U}_J(1/2J_{12})$, we had the expression
          %       $$\hat{U}_J(1/2J_{12})=\exp\left(-i\frac{\pi}{4}\hat{\sigma}_z^{(1)}\otimes\hat{\sigma}_z^{(2)}\right).$$
          %       This means that
          %       $$
          %         \hat{U}_J(1/2J_{12})\hat{R}_y^{(2)}(\pi/2)
          %         =\exp\left(-i\frac{\pi}{4}\hat{\sigma}_z^{(1)}\otimes\hat{\sigma}_z^{(2)}\right)\exp\left(-i\frac{\pi}{4}\hat{I}^{(1)}\otimes\hat{\sigma}_y^{(2)}\right)
          %         =\exp\left(-i\frac{\pi}{4}\left[\hat{\sigma}_z^{(1)}\otimes\hat{\sigma}_z^{(2)}+\hat{I}^{(1)}\otimes\hat{\sigma}_y^{(2)}\right]\right)
          %       $$
          % \item \begin{align*}
          %         \hat{R}_z^{(2)}(-\pi/2)\hat{U}_J(1/2J_{12})\hat{R}_y^{(2)}(\pi/2) &
          %         =\exp\left(i\frac{\pi}{4}\hat{I}^{(1)}\otimes\hat{\sigma}_z^{(2)}\right)\exp\left(-i\frac{\pi}{4}\left[\hat{\sigma}_z^{(1)}\otimes\hat{\sigma}_z^{(2)}+\hat{I}^{(1)}\otimes\hat{\sigma}_y^{(2)}\right]\right)                                        \\
          %                                                                           & =\exp\left(-i\frac{\pi}{4}\left[-\hat{I}^{(1)}\otimes\hat{\sigma}_z^{(2)}+\hat{\sigma}_z^{(1)}\otimes\hat{\sigma}_z^{(2)}+\hat{I}^{(1)}\otimes\hat{\sigma}_y^{(2)}\right]\right) \\
          %                                                                           & =\exp\left(-i\frac{\pi}{4}\left[(\hat{\sigma}_z^{(1)}-\hat{I}^{(1)})\otimes\hat{\sigma}_z^{(2)}+\hat{I}^{(1)}\otimes\hat{\sigma}_y^{(2)}\right]\right)
          %       \end{align*}
          % \item \begin{align*}
          %         \hat{R}_y^{(2)}(-\pi/2)\hat{R}_z^{(2)}(-\pi/2)\hat{U}_J(1/2J_{12})\hat{R}_y^{(2)}(\pi/2) & =\exp\left(i\frac{\pi}{4}\hat{I}^{(1)}\otimes\hat{\sigma}_y^{(2)}\right)\exp\left(-i\frac{\pi}{4}\left[(\hat{\sigma}_z^{(1)}-\hat{I}^{(1)})\otimes\hat{\sigma}_z^{(2)}+\hat{I}^{(1)}\otimes\hat{\sigma}_y^{(2)}\right]\right) \\
          %                                                                                                  & =\exp\left(-i\frac{\pi}{4}\left[(\hat{\sigma}_z^{(1)}-\hat{I}^{(1)})\otimes\hat{\sigma}_z^{(2)}+\cancel{\hat{I}^{(1)}\otimes\hat{\sigma}_y^{(2)}}-\cancel{\hat{I}^{(1)}\otimes\hat{\sigma}_y^{(2)}}\right]\right)             \\
          %                                                                                                  & =\exp\left(-i\frac{\pi}{4}\left[(\hat{\sigma}_z^{(1)}-\hat{I}^{(1)})\otimes\hat{\sigma}_z^{(2)}\right]\right)
          %       \end{align*}

          % \item \begin{align*}
          %         \hat{R}_z^{(1)}(\pi/2)\hat{R}_y^{(2)}(-\pi/2)\hat{R}_z^{(2)}(-\pi/2)\hat{U}_J(1/2J_{12})\hat{R}_y^{(2)}(\pi/2) & =\exp\left(-i\frac{\pi}{4}\hat{\sigma}_z^{(1)}\otimes \hat{I}^{(2)}\right)\exp\left(-i\frac{\pi}{4}\left[(\hat{\sigma}_z^{(1)}-\hat{I}^{(1)})\otimes\hat{\sigma}_z^{(2)}\right]\right) \\
          %                                                                                                                        & =\exp\left(-i\frac{\pi}{4}\left[(\hat{\sigma}_z^{(1)}-\hat{I}^{(1)})\otimes\hat{\sigma}_z^{(2)}+\hat{\sigma}_z^{(1)}\otimes \hat{I}^{(2)}\right]\right)                                \\
          %                                                                                                                        & =\exp\left(-i\frac{\pi}{4}\left[\hat{\sigma}_z^{(1)}\otimes\hat{\sigma}_z^{(2)}+\hat{\sigma}_z^{(1)}\otimes \hat{I}^{(2)}-\hat{I}^{(1)}\otimes\hat{\sigma}_z^{(2)}\right]\right)       \\
          %                                                                                                                        & =\exp\left(-i\frac{\pi}{4}\left[\hat{\sigma}_z^{(1)}\otimes(\hat{\sigma}_z^{(2)}+\hat{I}^{(2)})-\hat{I}^{(1)}\otimes\hat{\sigma}_z^{(2)}\right]\right)                                 \\
          %                                                                                                                        & =\exp\left(-i\frac{\pi}{4}\hat{\sigma}_z^{(1)}\otimes(\hat{\sigma}_z^{(2)}+\hat{I}^{(2)})\right)\exp\left(i\frac{\pi}{4}\hat{I}^{(1)}\otimes\hat{\sigma}_z^{(2)}\right)
          %       \end{align*}
          %       Now let's look at the second term. If we use Euler's formula to write it as a sum of sin and cos and Taylor expand both we see that
          %       \begin{align*}
          %         \exp\left(i\frac{\pi}{4}\hat{I}^{(1)}\otimes\hat{\sigma}_z^{(2)}\right) & =\cos\left(\frac{\pi}{4}\hat{I}^{(1)}\otimes\hat{\sigma}_z^{(2)}\right)+i\sin\left(\frac{\pi}{4}\hat{I}^{(1)}\otimes\hat{\sigma}_z^{(2)}\right)                                                                                                                           \\
          %                                                                                 & =\sum_{n=0}^{\infty}(-1)^n\frac{\left(\frac{\pi}{4}\hat{I}^{(1)}\otimes\hat{\sigma}_z^{(2)}\right)^{2n}}{(2n)!}+i\sum_{n=0}^{\infty}(-1)^n\frac{\left(\frac{\pi}{4}\hat{I}^{(1)}\otimes\hat{\sigma}_z^{(2)}\right)^{2n+1}}{(2n+1)!}                                       \\
          %                                                                                 & =\sum_{n=0}^{\infty}(-1)^n\left(\frac{\pi}{4}\right)^{2n}\frac{\left(\hat{I}^{(1)}\otimes\hat{\sigma}_z^{(2)}\right)^{2n}}{(2n)!}+i\sum_{n=0}^{\infty}(-1)^n\left(\frac{\pi}{4}\right)^{2n+1}\frac{\left(\hat{I}^{(1)}\otimes\hat{\sigma}_z^{(2)}\right)^{2n+1}}{(2n+1)!} \\
          %       \end{align*}
          %       Now we use that tensor products obey $(A\otimes B)(C\otimes D)=AC\otimes BD$
          %       and recall that $\hat{\sigma}_z^{2n}=\hat{I}$ (and of course $\hat{I}^{k}=\hat{I}$) which breaks the sums down to
          %       \begin{align*}
          %          & =\sum_{n=0}^{\infty}(-1)^n\left(\frac{\pi}{4}\right)^{2n}\frac{\left(\hat{I}^{(1)}\otimes\hat{I}^{(2)}\right)}{(2n)!}+i\sum_{n=0}^{\infty}(-1)^n\left(\frac{\pi}{4}\right)^{2n+1}\frac{\hat{I}^{(1)}\otimes\hat{\sigma}_z^{(2)}}{(2n+1)!} \\
          %          & =\cos\left(\frac{\pi}{4}\right)\hat{I}^{(1)}\otimes\hat{I}^{(2)}+i\sin\left(\frac{\pi}{4}\right)\hat{I}^{(1)}\otimes\hat{\sigma}_z^{(2)}                                                                                                  \\
          %          & =\frac{1}{\sqrt{2}}\hat{I}^{(1)}\otimes\left[\hat{I}^{(2)}+\hat{\sigma}_z^{(2)}\right].
          %       \end{align*}
          %       The other exponentials,
          %       $$
          %         \exp\left(-i\frac{\pi}{4}\hat{\sigma}_z^{(1)}\otimes(\hat{\sigma}_z^{(2)}+\hat{I}^{(2)})\right)
          %         =\exp\left(-i\frac{\pi}{4}\hat{\sigma}_z^{(1)}\otimes\hat{\sigma}_z^{(2)}\right)\exp\left(-i\frac{\pi}{4}\hat{\sigma}_z^{(1)}\otimes\hat{I}^{(2)}\right)
          %       $$
          %       expand in much the same way giving us
          %       \begin{align*}
          %         \exp\left(-i\frac{\pi}{4}\hat{\sigma}_z^{(1)}\otimes\hat{\sigma}_z^{(2)}\right)\exp\left(-i\frac{\pi}{4}\hat{\sigma}_z^{(1)}\otimes\hat{I}^{(2)}\right) & =
          %         \left[\cos\left(\frac{\pi}{4}\right)\hat{I}^{(1)}\otimes\hat{I}^{(2)}-i\sin\left(\frac{\pi}{4}\right)\hat{\sigma}_z^{(1)}\otimes\hat{\sigma}_z^{(2)}
          %         \right]                                                                                                                                                                                                                                                                                                      \\
          %                                                                                                                                                                 & \quad\left[\cos\left(\frac{\pi}{4}\right)\hat{I}^{(1)}\otimes\hat{I}^{(2)}-i\sin\left(\frac{\pi}{4}\right)\hat{\sigma}_z^{(1)}\otimes\hat{I}^{(2)}
          %         \right]                                                                                                                                                                                                                                                                                                      \\
          %                                                                                                                                                                 & =
          %         \frac{1}{2}\left[\hat{I}^{(1)}\otimes\hat{I}^{(2)}-i\hat{\sigma}_z^{(1)}\otimes\hat{\sigma}_z^{(2)}
          %         \right]
          %         \left[\hat{I}^{(1)}\otimes\hat{I}^{(2)}-i\hat{\sigma}_z^{(1)}\otimes\hat{I}^{(2)}
          %         \right]                                                                                                                                                                                                                                                                                                      \\
          %       \end{align*}
          %       so the whole thing becomes
          %       \begin{align*}
          %         &=\frac{1}{2\sqrt{2}}\left[\hat{I}^{(1)}\otimes\hat{I}^{(2)}-i\hat{\sigma}_z^{(1)}\otimes\hat{\sigma}_z^{(2)}
          %         \right]
          %         \left[\hat{I}^{(1)}\otimes\hat{I}^{(2)}-i\hat{\sigma}_z^{(1)}\otimes\hat{I}^{(2)}
          %         \right]\left[\hat{I}^{(1)}\otimes\hat{I}^{(2)}+\hat{I}^{(1)}\otimes\hat{\sigma}_z^{(2)}\right]     
          %       \end{align*}






  \end{enumerate}
\end{soln}

\end{document}